\documentclass[11pt,oneside,openany]{book}
\usepackage[paperwidth=6.5in,paperheight=9.5in,inner=.72in,outer=.72in,top=.75in,bottom=.78in,headheight=14pt]{geometry}
\usepackage{fontspec}
\setmainfont{lmroman10-regular.otf}[BoldFont=lmroman10-bold.otf,ItalicFont=lmroman10-italic.otf,BoldItalicFont=lmroman10-bolditalic.otf,SmallCapsFont=lmromancaps10-regular.otf]
\usepackage[main=english]{babel}
\usepackage{amsmath,amssymb,graphicx,microtype,fancyhdr,needspace}
\usepackage[unicode,hidelinks]{hyperref}
\hypersetup{pdftitle={Carl Friedrich Gauss, Works, Volume VI: Astronomical Writings},pdfauthor={Carl Friedrich Gauss},pdfsubject={English translations of the astronomical writings collected in Werke, volume VI (Göttingen, 1874), with historical editorial matter and remarks by Ernst Schering},pdfkeywords={Carl Friedrich Gauss, astronomy, English translation, Werke VI, Ernst Schering}}
\setlength{\parindent}{1.2em}\setlength{\parskip}{0pt plus .3pt}
\setlength{\emergencystretch}{1.4em}
\pagestyle{fancy}\fancyhf{}
\fancyhead[L]{\small ELLIPTIC ELEMENTS OF PALLAS}
\fancyfoot[C]{\thepage}\renewcommand{\headrulewidth}{.25pt}
\newcommand{\SourceLeaf}[1]{\hypertarget{native-#1}{}\label{native-#1}}
\newcommand{\SourceGathering}[1]{\marginpar{\raggedright\scriptsize\bfseries #1}}
\newcommand{\Article}[1]{\par\medskip\Needspace{4\baselineskip}\pdfbookmark[1]{Art. #1}{palladis-art-#1}\centerline{\textbf{#1.}}\nobreak\smallskip}
\renewcommand{\thefootnote}{*)}
\newcommand{\HistoricalNote}[1]{\footnote{\csname HistoricalNoteText#1\endcsname}}
\newcommand{\BareRoot}[2][]{\mathord{{}^{#1}\!\surd}#2}
\widowpenalty=10000\clubpenalty=10000\displaywidowpenalty=10000
\raggedbottom

\newfontfamily\AstroSymbols{DejaVuSans.ttf}
\newcommand{\AscendingNode}{\text{\mdseries\AstroSymbols ☊}}
\newcommand{\MeanMotionGlyph}{\mathord{7}}

\newfontfamily\LatitudeSymbols{Liberation Serif}
\newcommand{\ApproxBeta}{\text{\mdseries\LatitudeSymbols β}}
\newcommand{\ObservedBeta}{\text{\mdseries\LatitudeSymbols ϐ}}

\newcommand{\DescendingNode}{\text{\mdseries\AstroSymbols ☋}}
